\documentclass[10pt,a4paper]{amsart}
\usepackage[margin=19mm]{geometry}
\usepackage{amsmath,amssymb,mathtools}
\usepackage[scr=rsfs,cal=euler]{mathalfa}
\usepackage{xfrac}
\usepackage[unicode,hidelinks,bookmarks=false]{hyperref}
\newcommand{\ie}{i.e.\ }
\newcommand{\cf}{cf.\ }
\newcommand{\resp}{respectively\ }
\newcommand{\Subsection}{\subsection}
\providecommand{\nobreakdash}{\nobreak\hbox{-}}
\setlength{\emergencystretch}{2em}
\multlinegap0pt
\usepackage{amssymb}
\newtheorem{guess}{Guess}
\newtheorem{theorem}{Theorem}
\newtheorem{proposition}{Proposition}
\title{Quasicrystal Scattering and the Riemann Hypothesis}
\author{Michael Shaughnessy}
\date{Source: arXiv:2410.03673v12 (CC BY 4.0)}
\begin{document}
\maketitle

\noindent\emph{Source and licence.} Quasicrystal Scattering and the Riemann Hypothesis. Version arXiv:2410.03673v12 (CC BY 4.0), distributed by arXiv under the Creative Commons Attribution 4.0 International licence, http://creativecommons.org/licenses/by/4.0/, as recorded in the arXiv licence metadata for this exact version and shown on the abstract page https://arxiv.org/abs/2410.03673v12. Full licence notice: \texttt{licenses/arxiv-2410.03673-CC-BY-4.0.txt}.\par\smallskip

\noindent\textit{Editorial scope.} Counted unit: the article's theorem thm:RH (source0.tex lines 857--861). The article's complete original proof of it is delivered with it (source0.tex lines 863--914, one \texttt{proof} environment of 52 lines). No mathematics is written, repaired or re-derived: the statement and the proof are the article's own lines, in the article's own notation, and no retained line is substituted or re-flowed.  Retained uncounted as the source-local context the proof needs: lines 656--670; lines 736--747; lines 825--831.  Nothing was omitted from any retained span, so the package carries no omission record.  No retained line contains a citation, so the wrapper carries no bibliography and no bibliography entry.  No retained line contains a figure environment, an image-inclusion command or drawing code, so no image is delivered and the package declares no asset dependency.  Editorial formatting only, in addition: an AMS \texttt{amsart} preamble that restates the article's own declarations and macros used by the retained lines, namely the environments \texttt{guess}, \texttt{theorem}, \texttt{proposition} on separate counters, together with the standard packages those lines need.  The retained environments print in this document as Guess 1, Theorem 1, Proposition 1 in order of appearance, whereas the article prints the same items with the numbering of its own full document.  The complete native source of the article as distributed in the arXiv e-print for this exact version is bundled unchanged as \texttt{sources/166306/source0.tex}.
\begin{guess}[Dyson]
\label{guess:chi}
The prime distribution
$\chi(x) = \sum_{n=1}^\infty \delta(x - \ln p_n)$
is a Fourier quasicrystal:
\begin{equation}
  \hat\chi(k)
   \;=\; c_0 \;+\; \sum_{m=1}^\infty c_m\,
                    \delta\!\left(k - \frac{\gamma_m}{2\pi}\right),
  \label{eq:guess}
\end{equation}
for some constant $c_0$ and coefficients $\{c_m\}$, where
$\rho_m = \beta_m + i\gamma_m$ runs over the non-trivial zeros
of $\zeta(s)$.
\end{guess}

\begin{theorem}[Fourier self-duality on $\mathcal{S}'$]
\label{thm:selfdual}
For every $f \in \mathcal{S}'(\mathbb{R})$,
\begin{equation}
  \mathcal{F}\!\left[\mathcal{F}[f]\right](x)
   \;=\; f(-x).
  \label{eq:selfdual}
\end{equation}
This is a theorem of distribution theory, holding unconditionally
with no hypotheses on $f$ beyond
$f \in \mathcal{S}'(\mathbb{R})$.
\end{theorem}

\begin{proposition}[Bounded peak coefficients]
\label{prop:bounded}
Each coefficient $c_m$ in Guess~\ref{guess:chi} satisfies
$|c_m| < \infty$. Equivalently, the normalised peak amplitude
$p_L^{\beta_m - 1/2}/|\rho_m|$ stays bounded as
$L \to \infty$.
\end{proposition}

\begin{theorem}[Riemann Hypothesis]
\label{thm:RH}
Every non-trivial zero $\rho_m = \beta_m + i\gamma_m$ of
$\zeta(s)$ satisfies $\beta_m = 1/2$.
\end{theorem}

\begin{proof}
By Proposition~\ref{prop:bounded}, the normalised peak
amplitude $p_L^{\beta_m - 1/2}/|\rho_m|$ stays bounded as
$L = p_L \to \infty$. This is possible only if
\begin{equation}
  \beta_m \;\le\; 1/2.
  \label{eq:upper}
\end{equation}

Apply Theorem~\ref{thm:selfdual} (self-duality) to $\chi$:
\begin{equation}
  \mathcal{F}\!\left[\mathcal{F}[\chi]\right](x)
   \;=\; \chi(-x)
   \;=\; \sum_n \delta(x + \ln p_n).
\end{equation}
The right-hand side is a unit-weight sum of deltas on
$\{-\ln p_n\}$. Computing the left-hand side from
Guess~\ref{guess:chi} gives
\begin{equation}
  \mathcal{F}[\hat\chi](x)
   \;=\; c_0\,\delta(x) \;+\;
          \sum_m c_m\,e^{-i\gamma_m x},
\end{equation}
using $\mathcal{F}[\delta(k - k_m)](x) = e^{-2\pi i k_m x}$.
For this to equal the unit-weight sum of deltas
$\sum_n \delta(x + \ln p_n)$, the coefficients $c_m$ must be
\emph{exactly} those that make the Fourier series reproduce
this distribution --- in particular, they must be uniformly
non-vanishing, since the deltas $\delta(x + \ln p_n)$ all
carry weight $1$, not weight $0$.

This rules out the case $\beta_m < 1/2$: if any $\beta_m <
1/2$, then $c_m = -p_L^{\beta_m-1/2}\,e^{i\gamma_m \ln p_L}/
\rho_m \to 0$ as $L \to \infty$. The vanishing of $c_m$ would
remove the corresponding Fourier mode from the double-FT
reconstruction of $\chi(-\,\cdot\,)$, contradicting the
identity $\mathcal{F}^2[\chi] = \chi(-\,\cdot\,)$ which carries
all unit weights at $\{-\ln p_n\}$.

By the functional equation $\zeta(s) = 0 \Rightarrow
\zeta(1 - \bar s) = 0$, the alternative was already
inconsistent with~\eqref{eq:upper}: if $\beta_m < 1/2$ for
some $m$, then $1 - \beta_m > 1/2$ would arise for the
partner zero, violating~\eqref{eq:upper}.

The combined constraints --- (a) $\beta_m \le 1/2$ from the
one-dimensionality of $\hat\chi$ and the prime spacing
structure, and (b) $\beta_m \ge 1/2$ from the functional
equation applied to (a), or equivalently from the
non-vanishing of $c_m$ required by self-duality --- force
$\beta_m = 1/2$ for every non-trivial zero $\rho_m$.
\end{proof}

\end{document}
